\section{Giriş}

Doğal mantık, iki ifade arasında \emph{ileri gerektirme} ($a \sqsubset b$), \emph{geri gerektirme} ($a \sqsupset b$) ve \emph{eşdeğerlik} ilişkilerini ayırt eder \citep{maccartney2009natural}. Bu ilişkiler belirlenimci bir tersine çevirme cebirine uyar: öncül ile hipotezin yer değiştirmesi \FE$\leftrightarrow$\BE dönüşümünü yapar ve \EQ'yu değiştirmez. DiCo-NLI (SemEval-2027 Görev~2) bu cebiri bir değerlendirme hedefine dönüştürür. Her tersinir ifade çifti iki yönde de sınıflandırılır ve sistemler ağırlıklı F1'in yanında çift düzeyinde iki metrikle puanlanır: \emph{SoftCons} (iki tahmin tersine çevirme altında uyumlu mu?) ve \emph{HardCons} (ikisi de doğru mu?). Görev İngilizce, İspanyolca, Baskça ve karma dilli bir track içerir. İngilizce veri üzerindeki pilot çalışma, ince ayarlı kodlayıcıların bile çiftlerin \%16--20'sinde, yalnızca kod çözücü modellerin ise çok daha fazlasında tutarsız olduğunu bulmuştur \citep{apaolaza2026assessing}. Bu bulgu, özbağlanımlı dil modellerindeki tersine çevirme lanetini (reversal curse) hatırlatır \citep{berglund2024reversal}.

Bir sınıflandırıcının tersine çevirme cebirini kendiliğinden içselleştirmesini ummak yerine, değerlendirme verisinin her tersinir çiftin \emph{iki} yönünü de içerdiğini gözlemliyor ve bu yönleri \emph{ortak} olarak çözümlüyoruz. İki yönlü sonsal dağılımlar üzerinde, tersine çevirme eşlemesiyle birbirine bağlanmış bir uzmanlar çarpımı \citep{hinton2002training}, tasarım gereği tutarlı tahminler üretir. Bu yöntem ek eğitim gerektirmez ve aynı ilişkiye iki farklı açıdan bakan kanıtları birleştirir.

Katkılarımız şunlardır:
(i) dört track'in tamamında çok dilli kodlayıcı çalışması (XLM-R, mDeBERTa ve NLI ön eğitimli mDeBERTa; kendi verisiyle, sıfır-örnek ve birleşik eğitim);
(ii) \emph{ortak çift yönlü çözümleme}: HardCons'u eğitim maliyeti olmadan 72 değerlendirmenin tamamında artırır (ortalama +3.9 puan) ve eğitim sırasında uygulanan bir tutarlılık kaybından daha etkilidir;
(iii) ince ayarlı kodlayıcılarda yön tutarsızlığının küçük, \NO kararının ise asıl hata kaynağı olduğunu gösteren bir hata analizi;
(iv) iki veri kümesi gözlemi: SoftCons girdiden bağımsız sabit bir tahminciyle en yüksek değerine ulaşır ve \NO örnekleri yalnızca çift yapısından tanınabilir.

\section{Görev ve Veri}
\label{sec:task}

\paragraph{Etiketler ve track'ler.} Her örnek, \EQ, \FE, \BE ya da \NO (benzerlik veya karşıtlık gibi diğer tüm PhrasIS ilişkileri) olarak etiketlenmiş sıralı bir ifade çiftidir $(a,b)$. $\Rev(\EQ)=\EQ$, $\Rev(\FE)=\BE$ ve $\Rev(\BE)=\FE$'dir; \NO'nun tersi yoktur ve yalnızca F1'de kullanılır. Veri, görüntü altyazıları ve haber başlıklarından oluşturulmuş PhrasIS'in \citep{lopezgazpio2024phrasis} tersinir alt kümesinden türetilmiştir. İspanyolca ve Baskça sürümler makine çevirisiyle üretilmiş ve iki dilli uzmanlarca doğrulanmıştır. Track 1--3 tek dillidir (her biri 3{,}042 eğitim / 660 dev örneği). Track~4 dilleri karıştırır: her kaynak çift altı sıralı diller arası kombinasyonun hepsinde yer alır (en-es, es-en, \dots; 18{,}252 / 3{,}960 örnek). Bir en-es örneğinin tersi, ifadelerin yer değiştirdiği bir es-en örneğidir.

\paragraph{Metrikler.} $\mathcal{P}$ tersinir çiftler $(x, \bar{x})$ kümesi, $f$ de sistem olsun. Bu durumda
$\text{SoftCons} = \frac{1}{|\mathcal{P}|}\sum \mathbb{I}[f(x)\neq\NO \wedge f(\bar{x})=\Rev(f(x))]$ ve
$\text{HardCons} = \frac{1}{|\mathcal{P}|}\sum \mathbb{I}[f(x)=y \wedge f(\bar{x})=\bar{y}]$ olur; burada $y$ ve $\bar{y}$ altın etiketlerdir.
SoftCons altın etiketleri dikkate almaz, HardCons ise eşli doğruluktur.

\paragraph{İki gözlem.} \textbf{(O1)} SoftCons altın etiketten bağımsız olduğu için, her zaman \EQ tahmin eden ve girdiye hiç bakmayan bir sistem SoftCons $=\%100$ alır (Tablo~\ref{tab:main}). Bu nedenle SoftCons ancak F1 ve HardCons ile birlikte anlamlıdır.
\textbf{(O2)} Hem eğitim hem dev kümesinde her tersinir örnek iki yönde, her \NO örneği ise tam olarak tek yönde bulunur. Dolayısıyla bir örneğin yer değiştirmiş girdisinin dosyada bulunup bulunmaması, \NO'yu dev kümesinde \%100 kesinlik ve duyarlılıkla belirler. Test kümesi bu özelliği korursa, yapıyı bilen bir çözümleyici bundan yararlanabilir. Böyle bir çözümleyiciyi ayrı bir tanı satırı olarak (\textit{struct}) raporluyor, ana sistemimizde kullanmıyoruz.

\section{İlgili Çalışmalar}

\paragraph{Mantıksal tutarlılık.} Tutarlılık kısıtları NLI modellerine çekişmeli düzenlileştirme \citep{minervini2018adversarially} ve simetri ile geçişlilik için mantıktan türetilmiş eğitim kayıpları \citep{li2019logic} yoluyla eklenmiştir. Benzer fikirler soru yanıtlamaya da uygulanmıştır \citep{asai2020logic}. Tutarlılık kıyaslama kümeleri arasında BECEL \citep{jang2022becel} vardır; \citet{wu2025transitive} ise geçişli öz-tutarlılığı altın etiket olmadan değerlendirir. Bu çalışmalar tutarlılığı çoğunlukla eğitim sırasında zorlar; biz ise onu çıkarım anında kesin olarak zorluyor ve iki yaklaşımı karşılaştırıyoruz.
\paragraph{Tersine çevirme laneti.} ``A, B'dir'' biçiminde eğitilen özbağlanımlı dil modelleri ``B, A'dır'' çıkarımını yapamaz \citep{berglund2024reversal,lin2024delving}; bu durum ters yönlü eğitimi motive etmiştir \citep{golovneva2024reverse}. DiCo-NLI aynı simetriyi ifadeler arasındaki ilişkiler için sınar.
\paragraph{Artefaktlar.} NLI verilerinin etiketleme artefaktları \citep{gururangan2018annotation}, çevrilmiş değerlendirme verilerinin ise çeviri artefaktları \citep{artetxe2020translation} içerdiği bilinmektedir. O2 gözlemi de aynı türden yapısal bir artefakttır.

\section{Sistem}
\label{sec:method}

\paragraph{Yön duyarlı sınıflandırıcı.} Bir kodlayıcı çifti \texttt{[CLS]}\,$a$\,\texttt{[SEP]}\,$b$ olarak okur ve doğrusal bir katman dört etiket üzerinde $p_\theta(y \mid a,b)$ dağılımını verir. Her tersinir çiftin iki yönünü zaten içeren tüm eğitim örnekleri üzerinde çapraz entropiyle eğitiyoruz. XLM-R\textsubscript{base} \citep{conneau2020unsupervised}, mDeBERTa-v3\textsubscript{base} \citep{he2023debertav3} ve çok dilli NLI verisiyle \citep[XNLI ve makine çevirisiyle çoğaltılmış İngilizce NLI derlemleri;][]{conneau2018xnli,laurer2024less} ince ayar yapılmış bir mDeBERTa-v3\textsubscript{base} modelini karşılaştırıyoruz. Sonuncusu, cümle düzeyindeki NLI'dan ara görev aktarımını sınar; 3 sınıflı çıktı katmanı yeni bir 4 sınıflı katmanla değiştirilir.

\paragraph{Ortak çift yönlü çözümleme.} Her test örneği $(a,b)$ için girdisi tam olarak $(b,a)$ olan örneği yalnızca ifadeler ve dilleri üzerinden buluyoruz (\texttt{pair\_id} ya da başka bir kimlik kullanılmaz). Böyle bir örnek varsa
\begin{equation}
y^\star = \operatorname*{arg\,max}_{y}\, \log p(y|a,b) + \log p(\Rev(y)|b,a)
\label{eq:joint}
\end{equation}
seçilir ($\Rev(\NO)=\NO$). Bir çiftin iki örneğinin amaç fonksiyonları birbirinin aynadaki görüntüsüdür; böylece $(a,b)$ $y^\star$, $(b,a)$ ise $\Rev(y^\star)$ etiketini alır. Yer değiştirmiş karşılığı olmayan örneklerde doğrudan argmax kullanılır. Dev kümesinde skorlar \texttt{pair\_id} ile gruplamayla birebir aynıdır. Denk.~\ref{eq:joint}, aynı ilişkiyi zıt taraflardan gören iki uzmanın çarpımıdır. Bu yöntem (i) ortak olarak \NO etiketlenmeyen her çiftte SoftCons $=1$'i garanti eder, (ii) daha emin olan yönün kanıtı diğerini düzeltebildiği için test anı veri artırımı gibi çalışır ve (iii) modelden bağımsızdır, yeniden eğitim gerektirmez. Tanı amaçlı \textit{struct} varyantı ayrıca yer değiştirmiş karşılığı olmayan tüm örneklere \NO atar ve \NO'yu Denk.~\ref{eq:joint}'in dışında tutar (O2).

\paragraph{Tutarlılıkla düzenlileştirilmiş eğitim.} Eğitim tarafındaki karşılık olarak, \citet{li2019logic} yaklaşımına benzer biçimde $p_\theta(\cdot\mid a,b)$ ile ters çevrilmiş dağılım $\Rev\circ p_\theta(\cdot\mid b,a)$ arasına simetrik bir KL terimi ekliyoruz: $\mathcal{L} = \mathcal{L}_{\text{CE}} + \lambda\,\tfrac{1}{2}\big[\mathrm{KL}(p_{ab}\Vert \Rev\,p_{ba}) + \mathrm{KL}(\Rev\,p_{ba}\Vert p_{ab})\big]$. Bir çiftin iki yönü her zaman aynı mini yığına konur.

\paragraph{Diller arası aktarım.} Her track için ayrı bir modelin (\emph{kendi verisi}) yanında, yalnızca İngilizceyle eğitip modeli diğer track'lere sıfır-örnek uyguluyoruz (\emph{yalnız EN}) ve dört eğitim kümesinin birleşimi üzerinde tek bir model eğitiyoruz (\emph{birleşik}).

\section{Deney Düzeni}
\label{sec:setup}

Resmi test kümesi ancak Ocak 2027'de yayımlanacağı için açık dev kümesini test kümesi olarak kullanıyor ve model seçiminde hiç kullanmıyoruz. Bunun yerine \emph{eğitim} kaynak çiftlerinin \%10'unu doğrulama dilimi olarak ayırıyoruz (her track'te aynı çiftler; böylece diller arası sızıntı olmaz). Bu dilimde öğrenme oranını $\{2,5\}\cdot10^{-5}$ arasından (EN üzerinde) ve en iyi epoch'u (tek dilli eğitimde 10, karma eğitimde 4 epoch içinden) seçiyoruz. Seçim her zaman ayrı çözümlemeyle yapılır, dolayısıyla çözümleyiciden bağımsızdır. AdamW, 32'lik yığın ve \%10 ısınma kullanıyor, dizüstü bilgisayarın 18\,GB belleğine sığmak için kelime gömmelerini donduruyoruz (Ek~\ref{app:hparams}). Her yapılandırma üç tohumla çalıştırılır (Apple M3 Pro'da tek dilli koşu başına 3.5--6 dakika). Tüm skorlar resmi puanlayıcıdan alınmıştır.

\paragraph{Kaynaklar.} Yalnızca açık kodlayıcı ağırlıklarını ve resmi eğitim verisini kullanıyoruz.\footnote{\url{https://inigolopezgazpio.net/SemEval-2027-Task-2-DiCo-NLI/}} Tek harici veri, mDeBERTa-NLI'ın eğitildiği NLI derlemleridir (XNLI, çevrilmiş İngilizce NLI). PhrasIS kullanılmamıştır.

\section{Sonuçlar}
\label{sec:results}

\begin{table*}[t]
\centering\small
\setlength{\tabcolsep}{4pt}
\resizebox{\textwidth}{!}{%
\begin{tabular}{lccccccccccccccc}
\toprule
 & \multicolumn{3}{c}{EN} & \multicolumn{3}{c}{ES} & \multicolumn{3}{c}{EU} & \multicolumn{3}{c}{Karma} & \multicolumn{3}{c}{Ortalama} \\
\cmidrule(lr){2-4} \cmidrule(lr){5-7} \cmidrule(lr){8-10} \cmidrule(lr){11-13} \cmidrule(lr){14-16}
Sistem & F1 & Soft & Hard & F1 & Soft & Hard & F1 & Soft & Hard & F1 & Soft & Hard & F1 & Soft & Hard \\
\midrule
Rastgele & 25.3 & 18.1 & 5.5 & 25.3 & 18.1 & 5.5 & 25.3 & 18.1 & 5.5 & 25.1 & 18.4 & 6.2 & 25.2 & 18.1 & 5.7 \\
Sabit \EQ & 11.0 & 100.0 & 31.4 & 11.0 & 100.0 & 31.4 & 11.0 & 100.0 & 31.4 & 11.0 & 100.0 & 31.4 & 11.0 & 100.0 & 31.4 \\
\midrule
XLM-R & 69.1 & 82.7 & 72.1 & 54.2 & 66.1 & 51.9 & 61.1 & 80.5 & 63.7 & 63.4 & 79.3 & 64.9 & 62.0 & 77.1 & 63.1 \\
\quad + ortak çözümleme & 69.6 & 93.9 & 77.3 & 58.6 & \textbf{93.5} & 68.4 & 61.0 & \textbf{94.3} & 70.0 & 63.8 & 88.4 & 69.7 & 63.3 & \textbf{92.5} & 71.3 \\
\midrule
mDeBERTa & 79.6 & 87.1 & 81.0 & 77.3 & 85.6 & 78.3 & 76.5 & 82.7 & 75.7 & 73.4 & 81.8 & 73.8 & 76.7 & 84.3 & 77.2 \\
\quad + ortak çözümleme & 80.3 & 91.8 & 84.5 & 77.9 & 89.3 & 81.1 & \textbf{76.7} & 87.5 & 78.8 & 73.5 & 86.2 & 76.6 & 77.1 & 88.7 & 80.2 \\
\midrule
mDeBERTa-NLI & 82.8 & 88.9 & 84.6 & \textbf{79.0} & 85.0 & 79.5 & 75.1 & 81.8 & 75.1 & 74.7 & 82.4 & 74.7 & 77.9 & 84.5 & 78.5 \\
\quad + ortak çözümleme & 83.0 & 94.2 & 87.5 & 78.7 & 90.4 & 82.2 & 75.6 & 86.8 & 78.5 & 74.8 & 88.4 & 78.3 & 78.0 & 89.9 & 81.6 \\
\midrule
mDeBERTa + $\mathcal{L}_\text{cons}$ & 80.7 & 90.0 & 83.9 & 76.5 & 86.9 & 78.6 & 75.1 & 85.4 & 77.4 & 72.4 & 85.7 & 75.4 & 76.2 & 87.0 & 78.8 \\
\quad + ortak çözümleme & 80.8 & 92.5 & 85.3 & 76.8 & 91.1 & 81.2 & 75.3 & 88.0 & \textbf{79.3} & 72.5 & 88.6 & 77.2 & 76.3 & 90.0 & 80.8 \\
\midrule
mDeBERTa-NLI, 3 tohum topluluğu & 83.7 & 89.9 & 86.6 & 78.8 & 87.4 & 80.5 & 76.4 & 82.7 & 77.3 & \textbf{76.0} & 83.9 & 76.6 & 78.7 & 86.0 & 80.2 \\
\quad + ortak çözümleme & \textbf{84.3} & \textbf{94.9} & \textbf{89.9} & 78.9 & 91.3 & \textbf{82.7} & 76.2 & 85.9 & 79.1 & 75.9 & \textbf{88.9} & \textbf{79.5} & \textbf{78.9} & 90.3 & \textbf{82.8} \\
\midrule
\textit{mDeBERTa + struct (O2)} & \textit{91.2} & \textit{100.0} & \textit{89.5} & \textit{90.2} & \textit{100.0} & \textit{88.2} & \textit{88.7} & \textit{100.0} & \textit{86.5} & \textit{87.9} & \textit{100.0} & \textit{85.6} & \textit{89.5} & \textit{100.0} & \textit{87.5} \\
\bottomrule
\end{tabular}}
\caption{Dev sonuçları (\%). Öğrenen sistemler: 3 tohumun ortalaması (standart sapmalar Ek~\ref{app:std}'de); \emph{3 tohum topluluğu} satırları ortalama değil, tek bir topluluğun (olasılık ortalaması) sonucudur. \emph{ortak}: ortak çift yönlü çözümleme (Denk.~\ref{eq:joint}). Kalın: en iyi öğrenen (tanı dışı) değer. Son satır O2'yi kullanır; yalnızca tanı amaçlıdır.}
\label{tab:main}
\end{table*}


\paragraph{Kodlayıcılar.} mDeBERTa, XLM-R'ı ortalamada 14.7 F1 puanı geride bırakır (Tablo~\ref{tab:main}). XLM-R bu küçük kümelerde kararsızdır da: bir İspanyolca tohumu 0.41 doğrulama F1'ine çöktü. İngilizcede mDeBERTa (79.6 F1 / 81.0 HardCons), iki düzen birebir aynı olmasa da DeBERTa-v3 pilot temel sisteminin \citep[75 / 79;][]{apaolaza2026assessing} üzerindedir. Skorlar EN'den ES'e, EU'ya ve Karma'ya doğru düşer. Karma track'te en zor kombinasyon İspanyolca--Baskça'dır (en-es için 75.1'e karşı 71.6 F1).

\paragraph{Ortak çözümleme.} 72 değerlendirmenin tamamında (6 sistem $\times$ 3 tohum $\times$ 4 track) ortak çözümleme HardCons'u \emph{her seferinde} artırdı; ortalama artış $+3.9$ puandır. Kazanç mDeBERTa'da $+3.0$, XLM-R'da $+8.2$'dir: model ne kadar zayıf ve kararsızsa kazanç o kadar büyüktür. Ortak çözümleme SoftCons'u da her değerlendirmede ($+6.5$), F1'i ise hafifçe artırdı ($+0.5$; 72 değerlendirmenin 51'inde pozitif). Ortak çözümleme tohum topluluğuyla birbirini tamamlar. En iyi sistemimiz, NLI ön eğitimli mDeBERTa'nın ortak çözümlemeli 3 tohumlu topluluğudur ve ortalama 78.9 F1 ile 82.8 HardCons'a ulaşır. $y^\star=\NO$ olan çiftler tutarsız sayıldığı için SoftCons \%100'ün altında kalır.

\paragraph{NLI ön eğitimi.} NLI ile ince ayar yapılmış modelden başlamak İngilizcede ($+3.2$ F1) ve İspanyolcada ($+1.7$), Karma'da ise hafifçe ($+1.3$) fayda sağlar. XNLI'da bulunmayan bir dil olan Baskçada fayda sağlamaz ($-1.4$). Ortak çözümleme altında ortalama $+1.4$ HardCons ekler.

\paragraph{Eğitimde mi, çözümlemede mi tutarlılık?} Tutarlılık kaybı amaçlandığı gibi çalışır. Tutarsız çiftlerin payını düşürür (\%2.5$\rightarrow$\%1.7) ve ayrı çözümlemede HardCons'u $+1.6$, SoftCons'u $+2.7$ artırır; bedeli $-0.5$ F1'dir. Ancak ortak çözümleme tek başına bu HardCons kazancının yaklaşık iki katını verir ($+3.0$), ortak çözümlemenin altına kaybı eklemek ise yalnızca $+0.6$ getirir (80.2$\rightarrow$80.8). Dolayısıyla tutarlılığı çözümleme anında zorlamak, onu öğretmekten hem daha ucuz hem daha etkilidir.

\begin{table}[t]
\centering\small
\setlength{\tabcolsep}{2.5pt}
\begin{tabular}{@{}llccccc@{}}
\toprule
Model & Çöz. & Doğru & \NO & Tuts. & Yanlış & F1\textsubscript{\NO} \\
\midrule
XLM-R & ayrı & 63.1 & 9.9 & 13.0 & 14.0 & 33.1 \\
 & ortak & 71.3 & 7.5 & 0.0 & 21.2 & 33.0 \\
mDeBERTa & ayrı & 77.2 & 13.2 & 2.5 & 7.1 & 53.9 \\
 & ortak & 80.2 & 11.3 & 0.0 & 8.5 & 54.0 \\
\bottomrule
\end{tabular}
\caption{Altın tersinir çiftlerin sonucu (çiftlerin \%'si; dev, track ve tohum ortalaması): iki yön de doğru (Doğru), bir yöne \NO tahmin edilmiş, tutarsız (Tuts.) ya da tutarlı ama yanlış. F1\textsubscript{\NO}: \NO sınıfının örnek düzeyindeki F1'i.}
\label{tab:errors}
\end{table}


\paragraph{Hatalar nereden geliyor?} Tablo~\ref{tab:errors}, her altın tersinir çiftin sonucunu ayrıştırır. mDeBERTa'da çiftlerin yalnızca \%2.5'i yön açısından tutarsızdır; baskın hata, tersinir bir çift için \NO tahmin etmektir (\%13.2). \NO sınıfı yalnızca 54 F1'e ulaşır. Ortak çözümleme tüm tutarsızlıkları giderir ve bazı hatalı \NO tahminlerini de düzeltir (\%13.2$\rightarrow$\%11.3), çünkü emin olan bir yön diğerini geçersiz kılar. Tablo~\ref{tab:main}'deki tanı amaçlı \emph{struct} satırı, \NO kararının ne kadar önemli olduğunu ölçer: \NO'yu model yerine çift yapısıyla (O2) belirlemek ortalama F1'i 76.7'den 89.5'e çıkarır. Dolayısıyla açık veride görevin en zor kısmı, ilişkinin \emph{hangi yönde} geçerli olduğu değil, \emph{tersinir olup olmadığına} karar vermektir. Test verisi bu yapıyı gizlemedikçe, bu kısım aynı zamanda en kolay kestirme yoldur.

\begin{table}[t]
\centering\small
\setlength{\tabcolsep}{4pt}
\begin{tabular}{@{}lcccc@{}}
\toprule
Eğitim verisi & EN & ES & EU & Karma \\
\midrule
Kendi verisi & 84.5 & 81.1 & 78.8 & 76.6 \\
Yalnız EN & 84.0 & 80.7 & 68.5 & 54.4 \\
Birleşik & 84.8 & 83.6 & 81.5 & 77.7 \\
\bottomrule
\end{tabular}
\caption{Ortak çözümlemeli mDeBERTa'nın eğitim verisine göre dev HardCons değeri (\%, 3 tohum): track'in kendi verisi, yalnızca Track~1 (diğerlerinde sıfır-örnek) ya da dört track'in birleşimi.}
\label{tab:transfer}
\end{table}


\paragraph{Diller arası aktarım.} Yalnızca İngilizceyle eğitilen bir model İspanyolcada HardCons'unu korur (kendi verisiyle 81.1'e karşı 80.7; Tablo~\ref{tab:transfer}), ancak 3.5 F1 puanı kaybeder. \NO'yu daha seyrek tahmin eder (duyarlılık \%59 yerine \%43); bu, \NO örneklerinde F1'e mal olur ama tersinir çiftleri korur. Baskçada 10, karma track'te 22 HardCons puanı kaybeder; yani diller arası çiftler diller arası denetim gerektirir. Dört eğitim kümesini birleştirmek en iyi tek model ayarıdır (kendi verisiyle 80.2'ye karşı ortalama 81.9 HardCons); en büyük kazanç Baskça ($+2.7$) ve İspanyolcada ($+2.5$).

\section{Çalışma Planı ve Görev Dağılımı}
\label{sec:plan}

Bölüm~\ref{sec:method}'teki tüm bileşenler uygulandı. Final raporu (28 Aralık) için planımız:
(i) NLI ön eğitimi ile birleşik eğitimi 3 tohumlu toplulukta birleştirmek ve \emph{her} örneğin yer değiştirmiş girdisini modelin kendisiyle puanlamak (ters yön dosyada olmasa da);
(ii) \NO logitini doğrulama diliminde kalibre etmek ve ters çevrilmiş \NO eğitim örneklerini (bunlar da \NO'dur) 0.5 kayıp ağırlığıyla eklemek;
(iii) $\lambda\in\{0.1, 0.3, 1\}$ ayarı;
(iv) nitel hata analizi; isteğe bağlı olarak küçük, açık bir LLM ile sıfır-örnek deney (ayrı ve ortak çözümleme).

\begin{table}[h]
\centering\small
\setlength{\tabcolsep}{3pt}
\begin{tabular}{@{}lp{0.74\columnwidth}@{}}
\toprule
Üye & Kod / analiz; plan maddesi; yazdığı bölümler \\
\midrule
Küçük & Puanlayıcı entegrasyonu, temel sistemler, hata analizi (Tablo~\ref{tab:errors}); (iv); \S1, \S6 \\
Şahin & Eğitim altyapısı, ortak çözümleme, deneyler; (i); \S4, 5, 7, 8 \\
Okudurlar & Kodlayıcılar, NLI aktarımı, tutarlılık kaybı; (iii); \S3, \S6 \\
Tumurkhuu & Veri artefaktları (O1/O2), diller arası aktarım (Tablo~\ref{tab:transfer}); (ii); \S2, \S6 \\
\bottomrule
\end{tabular}
\caption{Görev dağılımı (taslağı herkes okur).}
\label{tab:plan}
\end{table}


\section{Sonuç}

DiCo-NLI'daki yön tutarlılığını bir çözümleme problemi olarak ele aldık. Ortak çözümleme maliyetsiz ve modelden bağımsızdır; eşli doğruluğu her değerlendirmede artırdı ve tutarlılık kaybından daha etkili oldu. Geriye kalan zorluk \emph{tersinir olmayan} ilişkileri tanımaktır. Açık veride bu karar çift yapısıyla atlanabilir; test verisi bu yapıyı gizlemelidir (ör.\ \NO örneklerini de iki yönde sunarak).

\section*{Sınırlılıklar}
Resmi test kümesi henüz yayımlanmadığı için tüm sayılar açık dev kümesi üzerindedir; model seçimi eğitim verisinden ayrılmış bir dilimle yapılmıştır. Ortak çözümleme şu an ters yönü yalnızca değerlendirme verisinde bulunduğunda kullanır; bu DiCo-NLI'da tasarım gereği geçerlidir, ancak tek yönlü NLI'da geçerli değildir. Yer değiştirmiş girdiyi modelin kendisiyle puanlamak (planlanan) bu kısıtı kaldırır. İspanyolca ve Baskça veriler çeviri olduğundan, çeviri artefaktları diller arası sonuçları şişirebilir. Hesaplama kısıtları nedeniyle (tek bir dizüstü GPU) kelime gömmeleri dondurulmuş temel boyutlu kodlayıcılar ve yapılandırma başına üç tohum kullandık.
